\subsection{Comparación global de estrategias}

La lectura conjunta de Grover y QFT permite identificar dos tendencias generales dentro de la
campaña experimental realizada. La primera es que la compresión resulta útil solo cuando el
estado conserva suficiente regularidad estructural. La segunda es que, dentro de ese régimen
favorable, Blosc ofrece la relación más equilibrada entre ahorro de memoria y sobrecarga
temporal, mientras que ZFP solo resulta competitivo en escenarios específicos y a costa de
introducir un error numérico pequeño pero no nulo.

\begin{table}[H]
    \centering
    \small
    \caption{Síntesis comparativa frente al almacenamiento no comprimido}
    \label{tab:summary-results}
    \begin{tabular}{>{\raggedright\arraybackslash}p{3.3cm}>{\centering\arraybackslash}p{1.8cm}>{\centering\arraybackslash}p{1.8cm}>{\centering\arraybackslash}p{1.8cm}>{\centering\arraybackslash}p{1.8cm}>{\centering\arraybackslash}p{2cm}}
        \hline
        \textbf{Carga de trabajo} & \textbf{Blosc: ahorro (\%)} & \textbf{Blosc: tiempo (x)} & \textbf{ZFP: ahorro (\%)} & \textbf{ZFP: tiempo (x)} & \textbf{ZFP: error} \\
        \hline
        Grover, objetivo único & 77.8 a 93.5 & 4.22 a 4.85 & 36.4 a 45.7 & 32.27 a 37.78 & $8.04\times10^{-15}$ a $1.61\times10^{-14}$ \\
        Grover, multiobjetivo & 77.8 a 93.4 & 4.17 a 4.86 & 36.7 a 45.7 & 33.12 a 38.13 & $8.04\times10^{-15}$ a $1.61\times10^{-14}$ \\
        QFT, superposición periódica & 51.4 a 75.3 & 2.54 a 3.48 & 74.2 a 83.2 & 4.41 a 5.75 & $3.79\times10^{-11}$ \\
        QFT, alta entropía & -49.0 a -12.1 & 10.24 a 12.52 & -55.5 a -2.1 & 85.65 a 103.01 & $2.82\times10^{-11}$ a $3.79\times10^{-11}$ \\
        \hline
    \end{tabular}
\end{table}

La evidencia experimental muestra que Grover constituye un escenario ampliamente favorable
para la compresión sin pérdida. En este circuito, Blosc ahorra más de 77\% de memoria en todos
los tamaños y supera 93\% en 24 qubits, con un costo temporal cercano a $4\times$. ZFP
introduce errores máximos absolutos del orden de $10^{-14}$, pero aun con esa tolerancia
reduce menos memoria y exige sobrecostos temporales superiores a $30\times$, por lo que su
ventaja práctica resulta claramente inferior dentro de esta campaña.

QFT obliga a matizar esa lectura. Cuando la entrada es periódica, ambos compresores siguen
siendo útiles: Blosc aporta una mejor relación entre ahorro y tiempo, mientras que ZFP logra
el mayor ahorro espacial a costa de errores del orden de $10^{-11}$ y de una penalización
temporal mayor. Sin embargo, cuando la entrada presenta alta entropía, la compresión fracasa
como estrategia general: el ahorro de memoria se vuelve negativo y el tiempo de ejecución crece
de forma muy pronunciada incluso en ZFP, pese a que el error admitido sigue siendo pequeño. En
este sentido, la evaluación delimita con claridad que la conveniencia de comprimir el vector de
estado no depende solo del circuito, sino de la estructura concreta del estado que debe
representarse.
